\documentclass[11pt]{article}

\usepackage[margin=1in]{geometry}
\usepackage{amsmath, amssymb}
\usepackage[T1]{fontenc}
\usepackage[utf8]{inputenc}
\usepackage{lmodern}
\usepackage{enumitem}
\usepackage{hyperref}

\setlist[itemize]{topsep=4pt, itemsep=2pt, parsep=0pt}

\title{6.8610 Spring 2026 Lecture 4 Notes: Transformers}
\author{Junru Ren with help from Codex}
\date{February 12, 2026}

\begin{document}
\maketitle

\section{Transformers}

\subsection{Motivation}
The lecture introduces the \textbf{Transformer} as the architecture that turned attention from a helpful seq2seq add-on into the main computational primitive. The goal is to produce strong contextualized representations for every token while avoiding the sequential bottlenecks of recurrent models.

\begin{itemize}
    \item In Lecture 3, attention improved seq2seq models by letting a decoder look back at all encoder states instead of compressing the source sentence into one vector.
    \item The next step is to ask whether attention itself can do most of the modeling work.
    \item We want each token to build a rich context-dependent representation by deciding how much to use information from every other token in the same sequence.
    \item We also want a model that is highly parallelizable, unlike RNNs whose hidden states must be computed sequentially.
    \item This leads to self-attention, Transformer encoders, Transformer decoders, and practical pretrained systems such as BERT and GPT-2.
\end{itemize}

\subsection{Key Concepts}
\textbf{Definition: Self-Attention}

Self-attention is a mechanism in which each token in a sequence computes how much it should attend to every token in that same sequence, including itself, and uses those weights to form a new representation.

\begin{itemize}
    \item \textbf{Contextualized embeddings:} the main goal is not merely to score the next word, but to produce a strong representation of each token after incorporating its context.
    \item \textbf{Query, key, and value vectors:} each token embedding is projected into three learned views.
    \begin{itemize}
        \item The \textbf{query} acts like a probe: what information is this token looking for?
        \item The \textbf{key} describes what a token offers for matching.
        \item The \textbf{value} contains the information that is actually passed forward once attention weights are decided.
    \end{itemize}
    \item \textbf{Scaled dot-product attention:} token-to-token similarity is computed by a dot product between a query and keys, scaled before a softmax.
    \item \textbf{Residual connections and layer normalization:} these stabilize deep networks and help preserve useful signal through many layers.
    \item \textbf{Positional encodings:} plain self-attention is permutation-invariant, so token order must be injected explicitly.
    \item \textbf{Multi-head attention:} multiple independent attention heads let the model capture different relationships in parallel.
    \item \textbf{Transformer encoder:} a block consisting of multi-head self-attention, a feed-forward network, residual connections, and layer normalization.
    \item \textbf{Transformer decoder:} similar to an encoder, but it uses masked self-attention and, in encoder-decoder settings, an additional cross-attention sublayer.
    \item \textbf{Pretraining and transfer:} Transformer encoders can be pretrained on unlabeled text and then fine-tuned for labeled downstream tasks.
\end{itemize}

\subsection{Core Models and Equations}
\textbf{Per-token projections.} For input embedding \(x_i\),
\[
q_i = W_Q x_i, \qquad k_i = W_K x_i, \qquad v_i = W_V x_i.
\]
The same three weight matrices are used for every token in the sequence.

\textbf{Attention scores.} For token \(i\) attending to token \(j\),
\[
s_{ij} = q_i^\top k_j.
\]

\textbf{Scaled attention weights.} Let \(d_k\) be the key dimension. Then
\[
a_{ij} = \operatorname{softmax}_j\!\left(\frac{q_i^\top k_j}{\sqrt{d_k}}\right)
= \frac{\exp\!\left(q_i^\top k_j / \sqrt{d_k}\right)}
{\sum_{\ell=1}^{n} \exp\!\left(q_i^\top k_\ell / \sqrt{d_k}\right)}.
\]
The lecture emphasized that dividing by \(\sqrt{d_k}\) keeps logits from becoming too large, which would otherwise make the softmax nearly one-hot and gradients too small.

\textbf{Output representation.}
\[
z_i = \sum_{j=1}^{n} a_{ij} v_j.
\]
For the slide example with the token ``The,'' the lecture walked through
\[
z_1 = 0.87\,v_1 + 0.12\,v_2 + 0.01\,v_3 + 0\,v_4.
\]

\textbf{Matrix form of self-attention.} If
\[
Q = X W_Q, \qquad K = X W_K, \qquad V = X W_V,
\]
then the entire attention head can be written as
\[
\operatorname{Attention}(Q,K,V)
= \operatorname{softmax}\!\left(\frac{QK^\top}{\sqrt{d_k}}\right)V.
\]

\textbf{Multi-head attention.} For head \(h\),
\[
\text{head}_h = \operatorname{Attention}(QW_Q^{(h)}, KW_K^{(h)}, VW_V^{(h)}).
\]
The heads are concatenated and linearly combined:
\[
\operatorname{MHA}(Q,K,V) = \operatorname{Concat}(\text{head}_1,\dots,\text{head}_H)W_O.
\]

\textbf{Feed-forward sublayer.} After attention, each token representation is passed through a position-wise feed-forward network:
\[
\operatorname{FFNN}(z_i) = W_2 \,\sigma(W_1 z_i + b_1) + b_2.
\]

\textbf{Residual connection.}
\[
f_{\text{residual}}(x) = f(x) + x.
\]
In a Transformer block, residual connections are applied around both the attention sublayer and the feed-forward sublayer.

\textbf{Layer normalization.} For token representation \(z_t = [z_{t,1}, \dots, z_{t,d}]\),
\[
\mu_t = \operatorname{mean}(z_t), \qquad \sigma_t = \operatorname{std}(z_t),
\]
\[
\hat{z}_{t,j} = \frac{z_{t,j} - \mu_t}{\sigma_t + \epsilon},
\qquad
y_{t,j} = \gamma_j \hat{z}_{t,j} + \beta_j.
\]
The lecture noted both pre-norm and post-norm variants:
\[
r_t^{\text{pre}} = f(\operatorname{LayerNorm}(z_t)) + z_t,
\qquad
r_t^{\text{post}} = \operatorname{LayerNorm}(f(z_t) + z_t).
\]

\textbf{Positional encodings.} The original Transformer adds a positional vector \(p_i\) to each token embedding:
\[
x_i' = x_i + p_i.
\]
For sinusoidal positional encodings,
\[
\operatorname{PE}_{(\text{pos},\,2i)} =
\sin\!\left(\frac{\text{pos}}{10000^{2i/d_{\text{model}}}}\right),
\]
\[
\operatorname{PE}_{(\text{pos},\,2i+1)} =
\cos\!\left(\frac{\text{pos}}{10000^{2i/d_{\text{model}}}}\right).
\]

\textbf{Masked self-attention.} In a decoder, future positions are masked:
\[
\operatorname{MaskedAttention}(Q,K,V)
= \operatorname{softmax}\!\left(\frac{QK^\top + M}{\sqrt{d_k}}\right)V,
\]
where
\[
M_{ij} =
\begin{cases}
0 & \text{if } j \le i, \\
-\infty & \text{if } j > i.
\end{cases}
\]
This preserves autoregressive generation.

\textbf{BERT objectives.}
\begin{itemize}
    \item \textbf{Masked language modeling:}
    \[
    \max_\theta \sum_{t \in \mathcal{M}} \log p_\theta(x_t \mid x_{\setminus \mathcal{M}})
    \]
    where \(\mathcal{M}\) is the set of selected masked positions.
    \item \textbf{Next sentence prediction:}
    \[
    p_\theta(\texttt{isNext} \mid s_1, s_2).
    \]
\end{itemize}

\subsection{Algorithms / Architectures}
\textbf{1. Self-attention for one head}
\begin{itemize}
    \item Start with token embeddings \(x_1, \dots, x_n\).
    \item Project each token into a query, key, and value vector using shared matrices \(W_Q, W_K, W_V\).
    \item For a fixed token \(i\), compute dot-product scores \(q_i^\top k_j\) against every token \(j\).
    \item Scale the scores by \(\sqrt{d_k}\) and apply a softmax to obtain attention weights \(a_{ij}\).
    \item Form the new token representation \(z_i = \sum_j a_{ij} v_j\).
    \item Repeat for all tokens; these computations can be parallelized across positions.
\end{itemize}

\textbf{2. Transformer encoder block}
\begin{itemize}
    \item Add positional information to the input embeddings.
    \item Run multi-head self-attention to produce context-aware token representations.
    \item Apply a residual connection and layer normalization.
    \item Pass each token independently through a feed-forward network.
    \item Apply another residual connection and layer normalization.
    \item Stack several such encoder blocks to obtain stronger contextualized representations.
\end{itemize}

\textbf{3. Multi-head attention}
\begin{itemize}
    \item Duplicate the attention mechanism with multiple independent parameter sets.
    \item Let each head compute its own attention pattern and output representation.
    \item Concatenate all head outputs.
    \item Project the concatenated representation back to the model dimension.
\end{itemize}

\textbf{4. BERT pretraining}
\begin{itemize}
    \item Use only Transformer encoders.
    \item Select 15\% of tokens for the masked language-model objective.
    \item Of those selected tokens, replace 80\% with \texttt{[MASK]}, 10\% with a random token, and leave 10\% unchanged.
    \item Train the model to predict the original selected tokens.
    \item In the original lecture material, also train next sentence prediction by deciding whether sentence 2 truly follows sentence 1.
\end{itemize}

\textbf{5. BERT fine-tuning}
\begin{itemize}
    \item Keep the pretrained encoder stack.
    \item Replace the top prediction layer with a task-specific head.
    \item Feed downstream inputs using the same tokenization as in pretraining.
    \item Often freeze many BERT weights initially, then gradually unfreeze or fine-tune more of them.
\end{itemize}

\textbf{6. Transformer decoder}
\begin{itemize}
    \item Use masked self-attention so token \(i\) can only attend to positions \(\le i\).
    \item In encoder-decoder settings, add a cross-attention sublayer where decoder queries attend to encoder outputs.
    \item During training, use teacher forcing: provide gold previous target tokens one step at a time.
    \item During generation, condition only on the prefix generated so far.
\end{itemize}

\textbf{7. GPT-style decoding}
\begin{itemize}
    \item Use decoder-only Transformer blocks.
    \item Apply masked self-attention at every layer.
    \item Predict the next token autoregressively.
    \item Reuse the same token embedding and positional-encoding machinery across the stacked decoder blocks.
\end{itemize}

\subsection{Important Intuitions from the Professor}
\begin{itemize}
    \item \textbf{Queries, keys, and values play different roles.} The lecture stressed that values are not redundant: separating them gives the model more freedom, because the same representation does not have to serve both as a matching signal and as the content being passed forward.
    \item \textbf{The square-root scaling is about optimization, not just notation.} Without scaling, dot products grow with dimensionality, producing overly sharp softmax distributions and poor gradient behavior.
    \item \textbf{Self-attention is not a brand-new concept relative to Lecture 3.} It is closely related to seq2seq attention; the main difference is that the query, key, and value all now come from the same sequence.
    \item \textbf{Residual connections are a safety mechanism.} They preserve the original signal so a sublayer does not need to learn a perfect transformation immediately in order for information to survive.
    \item \textbf{Layer normalization stabilizes both activations and training.} The professor emphasized that raw values can become too large even before thinking about backpropagation, and that this instability makes learning sensitive to initialization.
    \item \textbf{Parallelizability is one of the main practical reasons Transformers were so successful.} Most of the work before the softmax can be done for all tokens at once, unlike RNNs.
    \item \textbf{Positional encodings should be understood by their design goals more than by memorizing one formula.} The important desiderata are order information, relative-position cues, and the ability to generalize beyond lengths seen during training.
    \item \textbf{Multi-head attention gives the model multiple chances to discover structure.} Different heads can specialize in different patterns, such as local syntax, long-range dependencies, or punctuation-sensitive behavior.
    \item \textbf{BERT is primarily for representation learning, not free-form generation.} The lecture repeatedly contrasted encoder-only models, which excel at contextualized embeddings and classification, with decoder-only models, which are used to generate text.
\end{itemize}

\subsection{Examples}
\begin{itemize}
    \item \textbf{Cartoon self-attention example:} for the sequence ``The brown dog ran,'' the lecture manually computed attention scores for ``The'' and then formed
    \[
    z_1 = 0.87\,v_1 + 0.12\,v_2 + 0.01\,v_3 + 0\,v_4.
    \]
    This illustrates that a token may attend mostly to itself, somewhat to a nearby modifier, and hardly at all to unrelated tokens.
    \item \textbf{Permutation-invariance warning:} without positional encodings, self-attention behaves too much like a bag of words. Reordering tokens does not by itself encode the correct sequence order.
    \item \textbf{Machine translation:} the lecture used English-to-Czech translation to motivate Transformer decoders, masked self-attention, and cross-attention.
    \item \textbf{BERT masked-LM training:} a sentence such as ``The brown dog ran'' can be turned into ``The \texttt{[MASK]} dog ran,'' and the encoder must recover the original token from context.
    \item \textbf{Sentence-pair classification with BERT:} the \texttt{[CLS]} token can be used as a summary representation for tasks with one label per input sequence or per sentence pair.
    \item \textbf{Token-level prediction with BERT:} downstream tasks such as tagging can use the contextualized representation at each output position rather than a single global representation.
    \item \textbf{Causal masking in GPT-style models:} when generating token \(t\), the model is allowed to look at tokens \(1,\dots,t-1\) but not at future tokens \(t+1,\dots,n\).
\end{itemize}

\subsection{Common Pitfalls / Limitations}
\begin{itemize}
    \item \textbf{No positionality by default:} raw self-attention is permutation-invariant and must be augmented with positional information.
    \item \textbf{Quadratic cost in sequence length:} standard self-attention requires interactions between every pair of positions, which scales as \(O(n^2)\).
    \item \textbf{Attention alone does not provide nonlinearity.} The feed-forward network remains important for expressive power.
    \item \textbf{Optimization instability:} without scaling, residual connections, and normalization, deep attention networks are much harder to train.
    \item \textbf{Encoder-only models cannot directly generate arbitrary sequences.} BERT learns excellent contextualized embeddings but is not itself a decoder-style language generator.
    \item \textbf{Fine-tuning can fail if preprocessing changes.} The lecture specifically warned that downstream inputs should use the same tokenization regime as pretraining.
    \item \textbf{Very long contexts are expensive.} The professor noted that large context windows require additional engineering tricks beyond the vanilla attention computation.
\end{itemize}

\subsection{Connections to Other Lectures}
\begin{itemize}
    \item This lecture builds directly on Lecture 3's seq2seq attention. Cross-attention in a Transformer decoder is conceptually the same operation as the encoder-decoder attention introduced there.
    \item It also resolves a central limitation from Lecture 3: RNNs process sequences sequentially, while self-attention gives every token direct access to every other token and is much easier to parallelize.
    \item Compared with Lecture 2's word2vec and bag-of-words representations, Transformers produce \textbf{contextualized} token representations rather than one static vector per word type.
    \item Compared with Lecture 1's general language-model framing, this lecture provides the architecture that ultimately became the dominant implementation for modern large language models.
    \item BERT and GPT-2 preview two important branches of later NLP systems: encoder-style pretrained models for understanding tasks and decoder-style models for autoregressive generation.
\end{itemize}

\subsection{Exam-Relevant Summary}
\begin{itemize}
    \item A self-attention head computes
    \[
    q_i = W_Q x_i,\quad k_i = W_K x_i,\quad v_i = W_V x_i,
    \]
    then uses scaled dot-product attention to produce \(z_i\).
    \item Core attention formula:
    \[
    \operatorname{Attention}(Q,K,V) = \operatorname{softmax}\!\left(\frac{QK^\top}{\sqrt{d_k}}\right)V.
    \]
    \item The \(\sqrt{d_k}\) factor prevents overly large logits and overly sharp softmax distributions.
    \item Values are separate from queries and keys because the model benefits from different representations for matching versus content passing.
    \item A Transformer encoder block consists of multi-head self-attention, residual connections, layer normalization, and a feed-forward network.
    \item Positional encodings are necessary because self-attention alone is permutation-invariant.
    \item Multi-head attention lets different heads capture different relations in parallel.
    \item Standard self-attention is \(O(n^2)\) in sequence length, but it is still highly practical because the computations are parallelizable.
    \item BERT is an encoder-only Transformer pretrained with masked language modeling and, in the original version, next sentence prediction.
    \item Original BERT masking rule: choose 15\% of tokens; of those, 80\% become \texttt{[MASK]}, 10\% become random tokens, and 10\% stay unchanged.
    \item BERT is typically fine-tuned by attaching a task-specific head while preserving the pretrained encoder.
    \item A Transformer decoder uses \textbf{masked self-attention}; an encoder-decoder Transformer also uses \textbf{cross-attention}.
    \item GPT-2 is decoder-only: it generates text autoregressively and does not use encoder-side cross-attention.
\end{itemize}

\end{document}
